% TOP_PROBLEM: {"id": 2302021, "title": "Research Problems in Function Theory — Problem 2.21", "category_id": 8, "category": {"id": 8, "name": "analysis", "display_name": "Analysis", "description": "Limits, continuity, calculus, and function theory.", "slug": "analysis", "order_index": 8, "created_at": "2026-07-31T15:26:25.671Z"}, "status": "open"}
% FIRST_POSTED: null
% ENTRY_KIND: statement_only
% Imported from: batch11.tex; UnsolvedMath ID 2302021
% Compile twice: lualatex 2302021.tex
\documentclass[11pt]{article}
\usepackage{amsmath,amssymb,mathtools}
\usepackage{fontspec}
\usepackage{unicode-math}
\setmainfont{Latin Modern Roman}
\setmathfont{Latin Modern Math}
\usepackage{xurl}
\usepackage[unicode,colorlinks=true,linkcolor=blue,urlcolor=blue,pdftitle={UnsolvedMath Problem 2302021}]{hyperref}
\newcommand{\sref}[2]{\hyperlink{source:#1:#2}{[S#2]}}
\newcommand{\cataloguescope}[1]{\paragraph{Mathematical statement.}#1}
\begin{document}
% UnsolvedMath ID 2302021; source catalogue code AMR-022-2021
\setcounter{section}{2302020}
\section{Research Problems in Function Theory — Problem 2.21}\label{2302021}
{\small\sffamily UnsolvedMath ID 2302021 · Analysis}
\subsection{Definitions and mathematical statement}

\paragraph{Shared notation.}$\mathbb N=\{1,2,\ldots\}$, and $\mathbb Z,\mathbb Q,\mathbb R,\mathbb C$ denote the integers, rationals, reals and complex numbers. Any other conventions are specified locally.


Let $f:\mathbb C\to\mathbb C$ be entire and set $f^{\circ1}=f$, $f^{\circ(n+1)}=f\circ f^{\circ n}$.
A point $z_0$ has exact period $n$ if $f^{\circ n}(z_0)=z_0$ and $f^{\circ k}(z_0)\ne z_0$ for $1\le k<n$.
The source calls such a point a ``fixed point of exact order $n$''.
An exact-period-$n$ point is repelling if $|(f^{\circ n})'(z_0)|>1$.

\paragraph{Statement recorded in the supplied source.}
Does every transcendental entire function have a repelling periodic point, that is, do there exist $n\ge1$ and $z_0$ of exact period $n$ with $|(f^{\circ n})'(z_0)|>1$? \sref{2302021}{2}
Here ``fixed points'' has the source's all-periods meaning; the question does not assert the existence of a repelling point of period one. Update 2.21 records the affirmative result first proved by Baker. \sref{2302021}{2}
\subsection{Short English statement}
Must a transcendental entire function have a periodic point that repels nearby points?
\subsection{Sources}
\begin{description}
\item[\hypertarget{source:2302021:1}{S1}] ULAM.AI and the UnsolvedMath Contributors, \emph{UnsolvedMath: A Curated Collection of Open Mathematics Problems}, record 2302021 (AMR-022-2021). CC BY 4.0. \url{https://huggingface.co/datasets/ulamai/UnsolvedMath}.
\item[\hypertarget{source:2302021:2}{S2}] W. K. Hayman and E. F. Lingham, \emph{Research Problems in Function Theory} (2018), Chapter 2, Problem 2.21 and Update 2.21. \url{https://arxiv.org/abs/1809.07200}.
\end{description}
\label{end:2302021}
\end{document}
